First of all, 
we want to thank all the reviewers for their valuable comments,
and we will revise our manuscript accordingly.
In what follows, we summarize our responses to the questions.
*** 
# 118A 
**Question A-1**: Analyze complexity between HE3DB and [73].  
**Answer A-1**: 
In short, our method requires $1+2 \cdot \lceil n/5\rceil$ bootstrapping operations
for the MSB extraction of an $n$-bit encrypted plaintext,
while [73] requires $1+2 \cdot \lceil n/3 \rceil$ 
(or $1+2 \cdot \lceil n/4 \rceil$ when the noise is fresh). 
Hence, when $n$ becomes large, 
our method enjoys a better slope than [73].
We do have different decomposition parameters, and
we use FFT instead of NTT, which can jointly result in performance 
differences. We find out that without AVX,
we can still be $10$x--$22$x faster than the
implementation of [73] (AVX2 actually doesn't
affect TFHEpp. TFHEpp becomes 3.5x slower if AVX is disabled).
We will include the above discussions in the revised manuscript.

*** 
# 118B
**Question B-1**: Achieving malicious security over FHE?  
**Answer B-1**: 
Achieving malicious security in an FHE-based approach 
requires additional cryptographic techniques 
such as ZKP or OPRF [a]. 
However, current methods result in relatively high computational 
costs, and can be slower than TEE or MPC-based approaches. 
We do believe that providing malicious security for FHE
is an important and promising direction of research.

[a] Hao Chen et al., "Labeled PSI from Fully Homomorphic Encryption with 
Malicious Security", CCS, 2018.

**Question B-2**: Access pattern leaked? Can MPC or TEE be better?  
**Answer B-2**: 
We would first like to kindly point out that, similar to HEDA [91], 
HE$^3$DB performs a linear scan of the entire database for each query,
which ensures access pattern privacy.
We agree that MPC and TEE-based approaches can also
hide access patterns [96, 16]. 
We will add more clarifications
on these approaches in Section 1.2 of the manuscript.

**Question B-3**: Why HEDA 80-bit and HE3DB 128-bit?  
**Answer B-3**: 
We are sorry that this is a typo. HEDA does provide 128-bit
security, and we will revise the manuscript.

*** 
# 118C

**Question C-1**: Alg 1. looks overcomplicated. Can't we simply ... to obtain $ct_4$?  
**Answer C-1**: 
As discussed in Section 1 of [31] and Figure 2 of [75],
to correctly bootstrap a TFHE ciphertext, 
we have to ensure that the MSB of the encrypted plaintext is zero.
Therefore, before applying PBS on $ct_1$, 
it is necessary to reset the MSB of the plaintext
encrypted in $ct_1$, which is basically why we need to first produce
$ct_4$.
We will include the above explanations in the manuscript.

**Question C-2**: Isn't possible ... to directly 
evaluate the function? The proposed method ... faster and more accurate?  
**Answer C-2**: 
Different from HMSB, the output of a high-precision absolute function 
depends on all bits of the encrypted plaintext.
Consequently, we cannot use the shift-and-subtract strategy in Alg. 2 
to remove the lower bits as noises like HMSB.
For homomorphic min/max and sorting functions, we show that we are faster
than the open-source implementations of [28], and are both faster
and more accurate than [75] (Please kindly refer to Fig. 6 and Fig. 7),
which are two of the most recent state-of-the-art methods.

**Question C-3**: The given LWEstoRLWE algorithm may be slower than
[25].  
**Answer C-3**: 
While [25] is an highly accurate way of converting LWEs to RLWE,
we find two limitations in applying [25] to DB.
First, it is not known how to change parameters ($n$ and $q$)
when converting the LWE ciphertexts to the RLWE ciphertext using [25].
Second, the runtime of HomMod is independent of the number of LWE ciphertexts
to be converted, while [25] grows linearly.
As a result, in encrypted DB where we need to change encryption parameters
for aggregation, HomMod is easier to adopt and also runs faster than [25] in 
our experiments.
We will add the above explanations in the revised manuscript.

**Comment C-4**: Range of the plaintext?  
**Answer C-4**: 
We encoded the plaintext as integers in the range $[-2^{k-1}, 2^{k-1})$ and 
adopt the two's complement representation where the MSB of the value indicates 
where the number is positive or negative (the sign).
In two's complement, the MSB of $a-b$ represents the sign bit of the 
subtracted result, which is 0 if $a\geq b$,
and 1 for if $a  <  b$. 

**Comment C-5**: Most solutions can be ... the parameter.   
**Answer C-5**: We agree that [28] does support 32-bit precision, and
many methods can support unbounded depth using standard FHE bootstrapping. 
We will correct the mistakes and revise the expressions
in Table 3 to better explain the practical benefits of our method 
when applied in DB.



Thank you for your question. We appreciate your 
additional comment because this was an issue that 
could not be fully discussed due to the 700-word limit
to the above clarification. 

First of all, we would like to clarify that, when we
compile [73] with the CMake flag "WITH_NATIVEOPT=ON", the speed improvement 
was at most 1.2x instead of 5x. 
We believe that, as you suggest (and also pointed out
in [73]), since NATIVEOPT is a machine-specific 
optimization flag, the speed improvement varies from
machine to machine.
In this work, we are running on a server-grade Intel Xeon 6226R
processor, while Table 2 of [73] performs the experiment on a 
consumer-grade Intel Core i7-9700 CPU (in general,
consumer-grade CPU runs faster than server-grade CPU in
single-core mode). We do agree that these are most
implementation variations, and plan to include more complexity
analysis in our manuscript as described below.

Next, we would like to clarify that, in Table 4 of our
manuscript, we report the runtime for 4 and 16-bit
homomorphic $\geq$ comparison. We note that this is different
from a 4-bit MSB extraction. Specifically, the comparison
between 4-bit numbers requires the MSB extraction of a
5-bit number, for that the subtraction between two 4-bit
numbers can result in a 5-bit number. For example, if we
wish to compare -7 (4-bit) and 6 (4-bit), the subtracted
result is -13 (5-bit). Therefore, in Table 4, we actually
use a 5-bit MSB extraction for 4-bit comparison and 17-bit
MSB extraction for 16-bit comparison. Thus,
the performance difference 
reported in Table 4 of our manuscript differs from that in Table 2 of [73].
We plan to modify Table 4 to better present the performance results reported in [73].

Finally, in terms of the algorithmic and complexity difference
with respect to [73], we believe that the one bit difference is 
subtle and reflects the fundamental difference between our algorithm
and that in [73]. Using the language of [73],
for a message $m$
encrypted in a ciphertext $({\mathbf{a}}, b)$, the lower
bits of $m$ is cleared through the HomFloor function, which set
the least $\log{p}=\log{(q/\alpha)}$ bits of $m$ to be zeroes (in fact, there is an
error between ${\mathbf{\rm Dec}}({\mathbf{c}}', d')$ and 
$\alpha r(m)$, but we assume the errors to be small).
Then, HomSign repeatedly calls HomFloor to clear all the lower bits
of $m$. In contrast, in our method, we multiply $({\mathbf{a}}, b)$
by a factor of $2^{k-1}$, and take out a $k$-bit segment of $m$ using TFHE PBS.
The difference here is essential: [73] takes out a $q$-bit 
segment of $m$ by directly taking $({\mathbf{a}}, b)\bmod{q}$,
while we take out a $k$-bit segment of $m$ by multiplying 
$({\mathbf{a}}, b)$ with $2^{k-1}$ and then reduce the result
by $Q'$ (e.g., the real ciphertext modulus that can be as large as $2^{54}$).
Without the multiplication between $({\mathbf{a}}, b)$ with $2^{k-1}$,
$p=q/\alpha$ can only be set to as large as $2^{4}$. The main
reason here is that $q$ can only be as
large as $2N=4096$ in most AP/GINX type bootstrapping procedures for
efficiency reasons,
and an $\alpha=2^{8}$ is required to separate $m$ from the LWE errors
$e$ (therefore $p=2^{12}/2^{8}=2^{4}$). Consequently, [73] can only take 
out a $4$-bit (i.e., $p=2^{4}$) 
segment of $m$ using a single call to HomFloor, while our $k$-bit
HMSB can use the full bootstrapping capability of TFHE to take out a 
$5$-bit (i.e., $k=5$) segment of $m$ (and you are correct that
our $k$ is also affected by the size of ciphertext modulus $Q'$, similar to [73]),
resulting in an extra bit being taken out in every bit clearing iteration
(and therefore the complexity difference).
This is also the reason
why we do not have to differentiate $Q$ and $Q'$, so we have less parameter
requirements. We believe that the above
analysis is largely independent of the secret distribution. However,
you are correct that smaller secrets will result in smaller noise growths,
and the exact parameters can vary in practice. 
When $N$ and $Q$ (and $Q'$) become large, we may see further differences
as the methods of determining $p$ in [73] and $k$ in our work both change.
We think that the above detailed discussion is very interesting and we plan
to include as much as possible in our manuscript 
(if we do not have enough space in the main body of texts,
we will include it in the appendix).